\newpage
\section{Infrastructure Graph \emojiinfratitle}\label{sec-infra}


Infrastructure graphs, composed of Points of Presence (PoPs) connected by physical edges, play a vital role across various social sectors such as delivering electricity/water/gas resources via power/water/gas networks~\cite{xing2022graph, yang2023graph, liao2021review}, supporting transportation through road networks~\cite{wang2020traffic, jiang2022graph}, and facilitating information transmission via computer networks~\cite{rusek2020routenet, zhang2023tfe}. Managing and understanding the complex physical relationships within these graphs is essential for improving infrastructure management in key areas like service delivery, function forecasting, and network optimization. In computer networks, for instance, the state of a logical traffic routing path is directly influenced by the underlying physical fiber links. Predicting critical metrics like traffic delay and jitter depends on assessing the status of all involved fiber links~\cite{ferriol2023routenet, rusek2020routenet}, as these metrics are highly interdependent. The intricate physical connections in infrastructure graphs create valuable opportunities for employing GraphRAG techniques to enhance infrastructure management and optimization.


\subsection{Data Sources and Graph Construction Methods}
Despite only a little existing literature on designing GraphRAG for infrastructure graphs, in the following, we summarize as well as potential graph construction potentially benefits infrastructure graph management:

\begin{itemize}[leftmargin=*]
    \item \textbf{Power Network}: Nodes (generators, demands) are buses, and edge weights model the lines, which depend on the impedance between two buses. Each node could have features such as voltage magnitude and angle, example such as IEEE-30/118 power systems~\cite{zimmerman2010matpower}.

    \item \textbf{Water Network}: Nodes are junctions (pipe connections or users) and sources (reservoirs and tanks), and edges are pipes with water flow.~\cite{xing2022graph}

    \item \textbf{Gas Network}~\cite{wang2024multiscale, yang2023graph}: nodes, edges represent natural gas pipelines (user load, diversion, input). 

    Correlation graph of pipeline parameters: to capture the strength of the interaction between parameters, person correlation is used to construct the correlation graph of parameters in the pipeline. 

    Connection graph of pipeline network. The connection relationship between different pipelines are from the topology of the pipeline network. The ratio of the flow is the weight.

    \item \textbf{Transportation Network}: road network (road intersections as the nodes and road connections as the edges) to encode spatial dependency, road traffic flow, and speed forecasting problems. Metro/Bus station problem, metro or bus station topology where stations are nodes and edges are metro lines or bus routes. Region-level problems, such as the regular or irregular regions, are used as the nodes in the graph. Spatial dependency between different regions can be extracted from land use purposes. Traffic nodes have time-series features~\cite{jiang2022graph}. 

    Road-level graphs: sensor graphs, road segment graphs, road intersection graphs, and road lane graphs. In Sensor graphs, each sensor is a node and the edges are road connections. Road segment graphs, road intersection graphs and road lane graphs where nodes are road segments, road intersections and road lanes. 

    Region-level graphs: irregular region graphs, regular region graphs, OD graphs. Nodes are regions of the city. Regular region graphs have grid-based partitioning (CNN used here). Irregular region graphs include all other partitioning approaches, road-based and zip code-based. In the OD graph, the nodes are original region - destination region pairs and their connections are defined by spatial neighborhood or other similarities.

    Station-level graphs: subway station graphs, bus station graphs, bike station graphs, railway station graphs, car-sharing station graphs, parking lot graphs, and parking block graphs. Natural links between stations are used to defined the edges, e.g., subway or railway lines, or the road networks. 

    \item \textbf{Computer Network}: A computer network system typically comprises two interconnected graph layers: the logical and the physical. The logical layer represents the flow of data, where traffic is routed through multiple intermediate routers before reaching its destination. In this layer, nodes correspond to routers, and edges represent the logical paths that data packets traverse. Conversely, the physical topology layer refers to the underlying infrastructure, with Points of Presence (PoPs) from providers such as Comcast and Verizon serving as nodes, and the physical fiber links between them forming the edges. Studies focused on the logical layer often model networks as graphs, where nodes represent devices (e.g., switches, routers) and edges denote links (e.g., fiber-optic cables) or traffic paths~\cite{ji2020traffic, zheng2019gcn}. This graph-based representation captures both the forwarding behavior of the network (i.e., traffic interactions) and its connectivity (i.e., topological behavior), offering valuable insights for network management. For instance, optical topology-aware traffic engineering has proven effective in mitigating issues like fiber cuts and flash crowds~\cite{nance2023improving}. For the phycial layer, many previous research has focused on inferring the complete as well as aligning the physical and logical alyer.
\end{itemize}

\subsection{Task}
\begin{itemize}
    \item \textbf{Power Networks}: Fault detection, time series prediction, power flow calculation, data generation, network maintenance and operation 

    \item \textbf{Water Networks}: Estimating the node water heads and flows in all pipes given water network layout, pipe characteristics, nodal demands, water levels at the reservoirs, and head measurements at limited locations.

    \item \textbf{Natural Gas Networks}: Energy Consumption Prediction.

    \item \textbf{Transportation networks}: flow prediction (road-level, region-level, station-level), traffic speed/time/congestion prediction (road-level and region-level problem), time of arrival prediction, ride-hailing demand, traffic accident, traffic anomaly, parking availability, transportation resilience, urban vehicle emission, railway delay, lane occupancy

    \item \textbf{Computer Networking System}: traffic classification, IoT detection, traffic forecasting (latency, delay and jitter prediction), Intrusion Detection, 
\end{itemize}

\subsection{Query and Query Processor}

\subsection{Retriever}

\subsection{Organizer}

\subsection{Generator}

\subsection{Paper}
\begin{itemize}
    \item RAG-based Explainable Prediction of Road Users Behaviors for Automated Driving using Knowledge Graphs and Large Language Models
    \item Tfe-gnn: A temporal fusion encoder using graph neural networks for fine-grained encrypted traffic classification
    \item One Train for Two Tasks: An Encrypted Traffic Classification Framework Using Supervised Contrastive Learning
    \item RouteNet: Leveraging Graph Neural Networks for network modeling and optimization in SDN
    \item Encrypted Network Traffic Classification Using a Geometric Learning Model
    \item Flow Topology-Based Graph Convolutional Network for Intrusion Detection in Label-Limited IoT Networks
    \item GCN-TC: Combining Trace Graph with Statistical Features for Network Traffic Classification
    \item NetBench: A Large-Scale and Comprehensive Network Traffic Benchmark Dataset for Foundation Models
    \item Lens: A Foundation Model for Network Traffic
    \item Interactive AI with Retrieval-Augmented Generation
for Next Generation Networking
    \item ET-BERT: A Contextualized Datagram Representation with Pre-training Transformers for Encrypted Traffic Classification
    \item Large Language Models for Networking: Applications, Enabling Techniques, and Challenges
    \item Large Generative AI Models for Telecom: The Next Big Thing?
    \item Large Language Models for Networking: Workflow, Advances and Challenges
    \item OPTIMAL POWER FLOW USING GRAPH NEURAL NETWORKS
\end{itemize}